\ifdefined\bxthree\else\def\bxthree{BX3}\fi
\ifdefined\purpose\else\def\purpose{\textit{Purpose Layer}}\fi
\ifdefined\bounds\else\def\bounds{\textit{Bounds Engine}}\fi
\ifdefined\fact\else\def\fact{\textit{Fact Layer}}\fi

\section{Limitations and Future Work}
\label{sec:limitations}

The Bailout Protocol provides a formal, architecturally enforced mechanism for mandatory human accountability in multi-agent systems. However, several conditions exist under which the protocol's guarantees hold with reduced efficacy, and several extensions are required to reach full deployment readiness. This section enumerates those conditions and outlines the corresponding future work.

\subsection{Human Response Latency}
\label{sec:human-latency}

The Bailout Protocol's upstream accountability guarantee depends on a human anchor who is reachable, notified, and capable of issuing a resolution within a clinically reasonable time window. In deployment contexts involving geographically distributed teams, multi-timezone operations, or personnel with limited connectivity, the HUMAN\_ACKNOWLEDGED state may persist for extended durations.

The protocol's timeout handling (Section~\ref{sec:timeout-handling}) prevents indefinite stalling by logging staleness warnings and issuing periodic reminders, but it does not reduce the latency of the human response itself. A trigger that sits in HUMAN\_ACKNOWLEDGED for 12 hours represents a system in an unresolved constrained state for that duration. For low-stakes consequential contexts, this latency may be acceptable. For time-sensitive operational domains (e.g., irrigation control during an acute weather event), the latency introduces material risk.

The current protocol treats human latency as an environmental parameter and does not attempt to optimize the routing of BailoutTriggers based on time-zone, availability, or proximity to the triggering event. Future work formalizes adaptive timeout and availability-aware routing (Section~\ref{sec:adaptive-timeout}).

\subsection{Adversarial Conditions}
\label{sec:adversarial-conditions}

The Bailout Protocol's formal guarantees assume that the system hierarchy and the Forensic Ledger are intact and trustworthy. Under adversarial conditions — specifically, a compromised Fact Layer or a tampered Forensic Ledger — the protocol's enforcement guarantees are materially reduced.

A compromised Fact Layer could theoretically suppress BailoutTrigger events, falsify state machine transitions, or fabricate RESOLVED entries without a corresponding human resolution. Similarly, a tampered Forensic Ledger could eliminate the forensic record of a bailout event, undermining the auditability guarantee.

The current design relies on the immutability of the Forensic Ledger's SHA-256 chain (Section~5) and the architectural isolation of the Fact Layer from the Bounds Engine, but it does not formally guarantee that these protections hold under active attack. Future work includes formal verification of the Fact Layer's tamper-evidence properties (Section~\ref{sec:formal-verification}) and the definition of an explicit threat model that specifies the adversarial conditions under which the protocol remains reliable.

Additionally, the protocol assumes that the Purpose Layer human anchor is a single identified individual. It does not currently model scenarios in which the human anchor themselves is adversarial — i.e., where the designated human intentionally circumvents the protocol to execute a reserved action without proper authorization. Handling this failure class requires external governance mechanisms (e.g., dual-signature requirements for high-stakes reserved actions) that are outside the scope of the current protocol but are addressed in the BX3 Framework's governance layer.

\subsection{Scalability of the Escalation Path}
\label{sec:scalability}

The Bailout Protocol's escalation algorithm (Algorithm~\ref{alg:escalation}) traverses the system hierarchy recursively until it reaches the Purpose Layer owner of the originating node. In a flat, wide hierarchy (many machine nodes, shallow depth), escalation is fast and the risk of hierarchy errors is low. In a deep, narrow hierarchy — common in large-scale enterprise deployments with multiple intermediate management layers — each bailout event must traverse every intermediate machine node.

The bypass mechanism ensures that no intermediate node attempts autonomous resolution, but it does not eliminate the latency added by each hop in the chain. In a 10-level hierarchy, the propagation chain $\phi$ records 10 transit events and the total escalation latency accumulates across each hop. While the latency per hop is small (milliseconds in a healthy network), it compounds under degraded network conditions and becomes non-trivial in systems where the hierarchy spans cloud regions or operates over intermittent connectivity.

The protocol's guarantees hold in terms of correctness — every trigger reaches its designated human anchor regardless of hierarchy depth — but the practical latency of that delivery in deep hierarchies is a deployment consideration. Future work proposes distributed human anchor pools (Section~\ref{sec:distributed-anchors}) as a structural mitigation for deep-hierarchy latency.

\subsection{Exploitation Risk: Strategic Trigger Suppression}
\label{sec:exploitation-risk}

The Bailout Protocol creates a hard block on all actuator commands from a Bounds Engine node when a trigger enters the BOUNDED state. This hard block is a strong deterrent against reckless behavior by the Bounds Engine, but it also creates an incentive structure that could be exploited by a sophisticated actor.

Specifically, a Bounds Engine node that wants to avoid the hard block could theoretically suppress the firing of a BailoutTrigger — declining to transition to BOUNDED even when a trigger condition is satisfied — and instead attempt an autonomous resolution. The Fact Layer's enforcement model validates proposed actions against Fact Layer enforcement rules (Section~\ref{sec:trigger-conditions}, T4), but it does not independently monitor whether the Bounds Engine is correctly identifying and firing trigger conditions. There is a gap between the enforcement of the protocol's state transitions and the completeness of the trigger detection itself.

The primary mitigation is that suppression of a legitimate BailoutTrigger is a forensic event: if the Bounds Engine suppresses a trigger and subsequently causes a Fact Layer violation, the Forensic Ledger records both the violation and the absence of a corresponding bailout event at the originating node. Post-hoc forensic analysis can detect the suppression pattern. However, this is a reactive mitigation, not a preventive one. The current protocol does not provide architectural prevention of trigger suppression.

Future work addresses this gap through formal verification of the trigger detection completeness (Section~\ref{sec:formal-verification}) and through the introduction of an independent trigger monitoring layer that cross-validates trigger detection without access to the originating node's reasoning.

\section{Future Work}
\label{sec:future-work}

The limitations described in Section~\ref{sec:limitations} motivate three primary directions for future research and engineering: formal verification of the protocol's core properties, adaptive timeout and availability-aware routing for the escalation algorithm, and distributed human anchor pools to address hierarchy depth and latency.

\subsection{Formal Verification of Core Properties}
\label{sec:formal-verification}

The Bailout Protocol's claims — that triggers always reach a human anchor, that no machine actor can autonomously resolve a trigger, and that the Forensic Ledger provides a complete and tamper-evident record — are currently supported by architectural argument and state machine specification. To make these claims deployment-ready for regulated or safety-critical contexts, they must be subjected to formal verification.

Formal verification of the Bailout Protocol proceeds along three axes:

\bigskip

\noindent\textbf{State machine correctness.} The transition function $\delta : \Sigma \times E \rightarrow \Sigma$ (Section~\ref{sec:transition-conditions}) must be verified to be total, deterministic, and complete: every state in $\Sigma$ and every event in $E$ has a defined transition, and no transition leads to an undefined state. Model checking tools (e.g., NuSMV, TLA+) applied to the Fact Layer implementation can verify these properties against the specification.

\bigskip

\noindent\textbf{Trigger completeness.} The set of trigger condition classes $\mathbb{K}$ (Section~\ref{sec:trigger-conditions}) must be verified to be complete with respect to the full space of authority-exceeded conditions in the system. This requires a formal characterization of the mandate resolution space and a proof that every boundary condition maps to at least one trigger class in $\mathbb{K}$. This is the verification most relevant to the exploitation risk described in Section~\ref{sec:exploitation-risk}.

\bigskip

\noindent\textbf{Forensic Ledger integrity.} The SHA-256 chaining mechanism of the Forensic Ledger (Section~5) must be verified to be append-only and tamper-evident: that any modification to a historical entry results in a detectable hash discontinuity, and that the chain cannot be reordered or truncated without detection. Cryptographic verification using the Bellare-Neven model for hash chains provides the formal foundation.

\bigskip

The BX3 Framework's formal verification program is documented separately in the BX3 Formal Methods Report (Bxthre3 Inc., 2026). The verification work for the Bailout Protocol is scoped as a continuation of that program.

\subsection{Adaptive Timeout and Availability-Aware Routing}
\label{sec:adaptive-timeout}

The current timeout constants $T_{\text{bound}} = 300$ seconds and $T_{\text{human}} = 7200$ seconds are static values set at system initialization. They represent a conservative compromise between allowing thorough human review and preventing indefinite stalling. In production deployments, a static timeout is a suboptimal solution: it is too aggressive for complex resolutions requiring domain expertise, and too passive for simple resolutions that a responsive human could complete in minutes.

Future work defines an adaptive timeout mechanism that adjusts $T_{\text{human}}$ dynamically based on:

\begin{enumerate}\setlength\itemsep{0.25em}
\item The complexity class of the triggering condition $\kappa \in \mathbb{K}$ (e.g., CONSTRAINT\_UNCOVERED events may require longer review than SCOPE\_EXCEEDED events).
\item The availability and current workload of the designated human anchor, as reported by the INBOX routing system.
\item Historical resolution times for comparable bailout events in the same operational domain.
\end{enumerate}

Additionally, the escalation algorithm (Algorithm~\ref{alg:escalation}) should be extended to support availability-aware routing: when multiple candidate human anchors are available (e.g., in a distributed team), the algorithm may select the anchor with the lowest current latency while maintaining the single-anchor guarantee. This requires a formal definition of the availability signal and a proof that availability-aware routing does not violate the no-cloning property of the BailoutTrigger.

The output of this future work is an adaptive routing extension to the escalation algorithm and a configurable timeout policy engine that respects domain-specific resolution requirements.

\subsection{Distributed Human Anchor Pools}
\label{sec:distributed-anchors}

The current protocol assigns a single designated human anchor to each Bounds Engine node: the Purpose Layer owner $\text{owner}(B)$. This single-anchor design is intentional and correctness-preserving — it ensures a single accountable decision-maker for each bailout event. However, in deployments where the hierarchy is deep (Section~\ref{sec:scalability}) or where the designated anchor may be unavailable for extended periods, the single-anchor model introduces latency and unavailability risks.

Future work defines a distributed human anchor pool mechanism that maintains the single-anchor guarantee while providing redundancy and reduced escalation latency:

\begin{enumerate}\setlength\itemsep{0.25em}
\item Each Purpose Layer owner defines a pool of $n \geq 1$ authorized human anchors, each of whom is capable of issuing resolutions for the owner's node.
\item The pool is ordered by declared availability and proximity to the triggering event's operational domain.
\item The escalation algorithm routes to the highest-ranked available anchor in the pool, rather than to a single pre-designated individual.
\item The anchor who issues the resolution is recorded in the forensic chain $\phi$, maintaining full traceability of the accountable decision-maker.
\item If the selected anchor is unreachable within $T_{\text{human}}$, the algorithm attempts the next anchor in the pool (with a proportional timeout extension), ensuring that unavailability of a single anchor does not stall the protocol.
\end{enumerate}

The distributed anchor pool preserves the single-anchor guarantee at the resolution level — exactly one human makes the decision — while eliminating the single point of failure in the routing path. The pool mechanism is architecturally analogous to redundant safety systems in critical infrastructure: the redundancy improves reliability without changing the fundamental nature of the safety guarantee.

Formal specification of the distributed anchor pool mechanism, including proofs of single-anchor preservation and availability-aware routing correctness, is scoped as future work to be published as an addendum to this paper.

\section{Conclusion}
\label{sec:conclusion}

The Bailout Protocol, as specified in this paper, provides the first complete, formally specified, architecturally enforced mechanism for mandatory human accountability in a multi-agent system operating under the BX3 Framework. Its core contribution is a direct response to the core failure mode identified in Section~\ref{sec:background}: that autonomous agentic systems operating under broad mandates will inevitably encounter conditions that exceed their resolution authority, and that without a mandatory architectural bypass, those conditions result in uncontrolled system behavior with no human accountability.

The protocol resolves this failure mode through four complementary design commitments: a deterministic state machine (Section~\ref{sec:bp-state-machine}) that has no autonomous resolution path, a complete trigger condition taxonomy (Section~\ref{sec:trigger-conditions}) that identifies all known authority-exceeded conditions, an escalation algorithm (Algorithm~\ref{alg:escalation}) that bypasses all machine actors and routes to the designated Purpose Layer human anchor with complete diagnostic context, and a forensic ledger (Section~5) that creates an immutable, auditable record of every bailout event and its resolution. Together, these commitments deliver the upstream accountability guarantee stated formally in Section~\ref{sec:formal-definitions}: that every unresolvable exception state propagates to a human accountability anchor, bypassing all intermediate machine actors, with complete diagnostic context, and that the system never fails downward into algorithmic chaos.

The protocol's guarantees do not require that the human anchor make the correct decision, that the resolution be timely, or that the system recover quickly. They require only that the system never autonomously resolve an authority-exceeded condition and that every bailout event produces a complete forensic record. This minimal but essential guarantee is what makes the Bailout Protocol a safety-critical architectural component rather than an operational workflow recommendation.

The Bailout Protocol is not an isolated contribution. It is the runtime enforcement layer of the BX3 Framework's accountability architecture, where the Purpose Layer's authored mandates meet the hard constraints of the Fact Layer. The protocol's correctness depends on the integrity of the Bounds Engine's constraint evaluation (Section~4), the completeness of the trigger condition taxonomy (Section~\ref{sec:trigger-conditions}), and the tamper-evidence properties of the Forensic Ledger (Section~5). These components together constitute the full accountable autonomy stack — from mandate authorship through constraint enforcement to human escalation.

More broadly, the Bailout Protocol demonstrates that human accountability in autonomous systems is not a governance policy or a ethical guideline, but a computable, specifiable, and formally verifiable architectural property. It can be embedded, enforced, and audited at runtime. This is the contribution that the BX3 Framework offers to the field: a concrete, implementable answer to the question of how to build agentic systems that remain under meaningful human control, even when operating at full autonomous scope.

Future work extends the protocol's deployment readiness through formal verification of its core properties (Section~\ref{sec:formal-verification}), adaptive timeout and availability-aware routing (Section~\ref{sec:adaptive-timeout}), and distributed human anchor pools (Section~\ref{sec:distributed-anchors}). These extensions preserve the protocol's core guarantees while addressing the practical limitations identified in production-scale and multi-timezone deployments.

\end{document}